\section{Conclusion}

This report has traced a coherent path from physical measurement to high-quality image reconstruction across two distinct but mathematically unified modalities. In both ISM and CT, the core challenge is the same: inverting a compact linear operator whose naïve inversion is unstable, and compensating for incomplete or noisy data through principled regularisation or learned priors.

The practical experiments on CT reconstruction illustrate the strength of machine learning techniques and the importance of domain-specific training. Starting from Filtered Back-Projection, which degrades quickly, through iterative methods with hand-crafted priors, such as Tikhonov and LASSO, we managed to get closer to a satisfactory image reconstruction. To improve the results, we tried to use a pretrained DnCNN denoiser as prior, which yielded worse results compared to the classical methods. We then tried to train the DnCNN denoiser on the dataset itself, thus using a denoiser trained on the specific task, which resulted in a significant improvement in the output..

Taken together, the results reinforce a central insight of modern computational imaging: the forward model provides the physics, but reconstruction quality ultimately depends on how well the prior captures the statistics of the target domain. 